\section{Limitations \& Future Work}
\label{sec:limitations}

This section reflects on the boundaries of the evidence presented in \cref{sec:results,sec:statistical_analysis,sec:cost_modelling,sec:safety,sec:experiments} and the corresponding artifacts generated by the implemented harness in \cref{sec:benchmark_harness}. We first articulate limitations that qualify the findings and then delineate concrete, system-aligned extensions that build directly on the same code paths, datasets, and telemetry.

\subsection{Limitations of the Current Study}
\label{subsec:limitations_study}

\paragraph{Scope and external validity of the corpus.} All results derive from a curated prompt bank designed for Socratic tutoring on engineering estimation (\cref{subsec:datasets_prompt_bank}). We executed 120 matched prompts across models (10 per category; \cref{sec:results}) using the standardized adapter interface (\cref{subsubsec:adapter_interface}). While this corpus was purpose-built to probe the evaluation axes in \cref{sec:evaluation_criteria}, its topical and stylistic coverage remains necessarily finite. The prompt bank emphasized estimation heuristics and reflective questioning rather than, for example, highly technical derivations, multi-document synthesis, or domain-specific notation. Consequently, generalization to other curricula or discourse styles should be interpreted cautiously.

\paragraph{Contextualization and controlled conditions.} The contextualization experiments (\cref{subsec:contextualization_test}) used teacher-provided context and two-turn interactions under controlled timing and concurrency (\cref{subsubsec:core_functionality}). These conditions establish clean contrasts (e.g., ephemeral user vs persistent system context) but do not capture the full variability of live classrooms: topic drift, student typos, or long-horizon plans. The binary inclusion measure and the continuous fidelity score (\cref{sec:statistical_analysis}) indicate internal consistency (cf. \cref{fig:context_inclusion_diff}), yet the construct validity for open-ended, long-run tutoring remains only indirectly supported.

\paragraph{Human scoring and inter-rater reliability.} Pedagogical quality and contextual fidelity were adjudicated from human ratings using the rubrics in \cref{sec:scoring_rubrics}. Despite rater calibration and blinded transcripts, human judgment introduces subjectivity. Inter-rater reliability fell in an acceptable but not perfect band: quadratic-weighted $\kappa=0.469$ (Socratic) and $\kappa=0.498$ (Context), with ICC(2,$k$)=0.73 and 0.75 respectively (see \cref{subsec:inter_rater_reliability,sec:results}). The $\kappa$ computation relies on ordinal discretization of rubric scores, which can compress subtle differences; conversely, ICC assumes approximate normality of aggregated scores. These modeling assumptions bound the precision of effect estimates in \cref{fig:effect_size_heatmap}.

\paragraph{Multiple comparisons and decision aggregation.} Pairwise contrasts controlled false discovery using Benjamini--Hochberg $q$-values and reported median paired differences (Hodges--Lehmann) together with nonparametric effect sizes (rank-biserial $r$ and Cliff's $\delta$; \cref{sec:statistical_analysis}). The global Friedman tests (\cref{subsec:comparative_analysis}) were decisive for both $Q$ and the context score, but the decision rule ultimately aggregates heterogeneous axes via the weighted sum in \cref{eq:overall_score} using \cref{tab:criteria_weights}. Any alternative governance weighting (\cref{subsec:decision_sensitivity}) would induce different utility frontiers; we partially mitigated this by reporting sensitivity analyses and effect sizes rather than point-rankings alone.

\paragraph{Operational measurements and load shape.} Latency was measured client-side as round-trip time through provider APIs under a fixed scheduler configuration (bounded concurrency, cooldowns; \cref{subsubsec:core_functionality}). The observed advantages for Gemini (cf. \cref{fig:cdf_latency,fig:latency_forest}) may vary under different network conditions, regions, or request shapes (streaming vs non-streaming). Our load tests targeted approximately 30 concurrent users (\cref{subsubsec:rate_limiting_results}); this probes throttling behavior (HTTP 429) and backoff efficacy but does not reach enterprise-scale saturation, leaving tail-behavior under heavy contention only partially characterized.

\paragraph{Safety and privacy evaluation.} Safety analyses used automated PII detection with the sanitization hook (\cref{subsec:sanitization_procedures,lst:sanitize_hook}). The PII detector registered minimal incidents in aggregate (\cref{tab:pii_detection_counts}), validating the pipeline's effectiveness. However, this does not constitute an exhaustive red-team; the detection emphasizes false negative risk in the presence of novel identifiers or code-mixed inputs. Comprehensive safety evaluation would require targeted adversarial testing.

\paragraph{Cost model and price drift.} The cost model computes per-call spend from measured tokens and contemporaneous provider price sheets (\cref{eq:per_call_cost,sec:cost_modelling}). Published prices, free-tier policies, and tokenization algorithms change over time; moreover, hosted open models can exhibit host-specific pricing and tokenization. While we persisted price tables with each run (\cref{subsubsec:data_persistence}), any absolute projections (e.g., per-1k interactions) should be read as contingent on the recorded run metadata.

\paragraph{Feature boundary of the harness.} The present harness standardizes text-only chat completions without external tool use or retrieval (\cref{subsubsec:adapter_interface}). As such, we did not evaluate function-calling APIs, RAG pipelines, or multimodal inputs. The conclusions therefore apply to purely conversational Socratic tutoring with ephemeral context injection (\cref{subsec:contextualization_test}) and controlled prompt templates (\cref{subsec:few_shot_prompts}).

\paragraph{Category stratification and representativeness.} Analyses stratified by instructional category (\cref{sec:results}) reveal patterned differences, but the strata reflect our curated taxonomy (\cref{subsec:datasets_prompt_bank}). Curricular variants (e.g., discipline-specific estimation norms) could shift relative strengths. The harness records rich metadata (\cref{subsubsec:data_logged,subsubsec:data_persistence}), enabling future subgroup or fairness analyses that were out of scope here.

\subsection{Future Experiments and Research Directions}
\label{subsec:future_directions}

The limitations above motivate concrete, incremental extensions that reuse the same harness interfaces, telemetry schema, and analysis code. We outline the highest-impact directions and the precise artifacts they would exercise.

\paragraph{Longitudinal, in-situ learning outcomes.} Building on the category- and prompt-level metrics in \cref{sec:results}, a classroom deployment can connect harness telemetry (latency, inclusion, hallucination proxies; \cref{subsubsec:data_logged}) to student learning outcomes via pre/post assessments. A randomized A/B design (Socratic chatbot vs platform-only baseline) would quantify effect sizes beyond rubric scores, using the statistical procedures in \cref{sec:statistical_analysis} and monitoring gates aligned to the SLOs of \cref{subsec:slos_monitoring}. The same SQLite schema (\cref{subsubsec:data_persistence}) supports reproducible education-facing dashboards.

\paragraph{Retrieval-augmented grounding.} The conversational-only boundary in \cref{subsubsec:adapter_interface} can be extended with a retrieval layer over approved course materials. Evaluations would track shifts in hallucination proxies and fidelity (\cref{sec:safety,sec:statistical_analysis}) and incorporate retrieval overhead in \cref{eq:per_call_cost}. Contrasts would reuse \cref{subsec:contextualization_test} to compare retrieved snippets in user vs system positions, updating \cref{fig:context_inclusion_diff} and the cost projections in \cref{sec:cost_modelling} accordingly.

\paragraph{Multi-turn and long-horizon tutoring.} We emphasized two-turn probes for clean comparability (\cref{sec:experiments}). Follow-on experiments should exercise longer sequences with evolving student states to interrogate planning, consistency, and recovery from earlier missteps. Outcomes would extend the rubric in \cref{sec:scoring_rubrics} with session-level constructs (e.g., misconception repair over multiple turns) while reusing latency/throughput measurement (\cref{fig:cdf_latency}) and reliability accounting (\cref{subsubsec:core_functionality}).

\paragraph{Robustness, fairness, and subgroup analysis.} With the harness metadata in \cref{subsubsec:data_logged}, we can stratify by linguistic variants, prompt perturbations, or accessibility accommodations to probe robustness. Fairness-oriented analyses would add protected-attribute surrogates to the SQLite store (\cref{subsubsec:data_persistence}) under the data retention policy of \cref{subsec:data_retention} and evaluate stability of $Q$ and inclusion across groups with the same nonparametric toolkit of \cref{sec:statistical_analysis}.

\paragraph{Routing policies under load and drift.} The hybrid routing sketch in \cref{subsec:hybrid_approach} can be operationalized as controlled policies (static thresholds vs learned bandits) evaluated under the load scenarios of \cref{subsubsec:rate_limiting_results}. The objective combines quality and cost via \cref{eq:overall_score,eq:routing_cost}, with acceptance tests drawn from \cref{subsec:risks_mitigations,subsec:slos_monitoring}. Telemetry-driven recalibration uses the same sensitivity framework as \cref{subsec:sensitivity_analysis} with policy-level CIs.

\paragraph{Open-model specialization and cost trade-offs.} To test whether a specialized open model can match the pedagogical headroom of the recommended configuration (\cref{sec:decision}), one can incorporate fine-tuned hosted Llama variants into the adapter layer (\cref{subsubsec:adapter_interface}). The evaluation would reuse the exact prompt bank (\cref{tab:prompt_bank,tab:system_prompts}), report the same inferential statistics (\cref{sec:statistical_analysis}), and extend cost projections with finetuning amortization in \cref{sec:cost_modelling}.

\paragraph{Security and red-teaming expansion.} Beyond PII detection (\cref{tab:pii_detection_counts}), targeted red-team prompts and safety taxonomies can be serialized into the same CSV/SQLite artifacts (\cref{subsubsec:data_persistence}) to quantify incident rates and mitigation efficacy. The sanitization hook (\cref{lst:sanitize_hook}) provides a natural intervention point for pattern-based filters; their impact can be measured through the acceptance criteria in \cref{subsec:risks_mitigations}.

Collectively, these directions preserve the end-to-end reproducibility guaranteed by \cref{sec:reproducibility_guide}: each extension is formulated as additional harness workloads, labels, or adapters, with all outputs entering the same analysis pipeline. In this way, the study's limitations delimit a concrete roadmap that incrementally broadens curricular coverage, operational stress, and safety assurance while retaining the comparability and auditability of the present results.

\subsection{Risks, Mitigations, and Acceptance Criteria}
\label{subsec:risks_mitigations}

Looking forward to deployment and operational monitoring, we identify key risks that could affect the decision in \cref{sec:decision} and specify mitigations grounded in the harness infrastructure.

\paragraph{Provider drift.} Silent model updates or pricing changes could alter standing. Mitigation: monthly re-benchmark on a 20\% sample using the harness workflow in \cref{sec:benchmark_harness} and compare to the stored confidence bands in \cref{tab:summary_scores}. Acceptance: no axis deviates beyond the prior 95\% CI without a compensating improvement elsewhere; otherwise escalate to routing recalibration (\cref{subsec:hybrid_approach}).

\paragraph{Category regressions.} Category-specific weaknesses (\cref{subsec:category_differentiation}) may affect curricular modules. Mitigation: category-aware routing and prompt audits; acceptance: maintain per-category medians within one IQR of baseline.

\paragraph{Operational spikes.} Latency/429 bursts under load (\cref{subsubsec:rate_limiting_results}). Mitigation: enforce cooldowns and backoff as implemented (\cref{subsubsec:core_functionality}); acceptance: P95 latency $\leq 1.2\times$ baseline and success rate $\geq 95\%$ during load tests.

\paragraph{Safety incidents.} Any PII persistence or hallucination surge. Mitigation: blocklist expansion and stricter templates (\cref{subsec:sanitization_procedures,subsec:few_shot_prompts}); acceptance: zero persisted PII (\cref{tab:pii_detection_counts}).

\subsection{Service Objectives and Monitoring Cadence}
\label{subsec:slos_monitoring}

To support operational deployment, we propose the following Service Level Objectives (SLOs), computed from harness-style telemetry (\cref{subsubsec:data_logged}):

\begin{itemize}
	\item \textbf{Pedagogy SLO:} Quarterly median $Q$ within the 95\% CI reported in \cref{tab:summary_scores}; alert at CI breach.
	\item \textbf{Latency SLO:} Weekly P95 latency within $1.2\times$ the baseline of \cref{fig:cdf_latency}; alert at breach for two consecutive weeks.
	\item \textbf{Reliability SLO:} Weekly first-try success rate $\geq$ 95\%.
	\item \textbf{Safety SLO:} Zero persisted PII per \cref{tab:pii_detection_counts}; hallucination rate within historical IQR.
	\item \textbf{Cost SLO:} Monthly spend within 10\% of the scenario projections in \cref{subsec:projected_scenarios} for the configured routing ratio.
\end{itemize}

Monitoring runs monthly (full) and weekly (canary subset). Dashboards read from the SQLite artifacts (\cref{subsubsec:data_persistence}). Incidents trigger the QA loop in \cref{subsec:integration_runbook}.
